% Part: model-theory
% Chapter: models-of-arithmetic
% Section: non-standard-models

\documentclass[../../../include/open-logic-section]{subfiles}

\begin{document}

\olfileid{mod}{mar}{nst}
\section{مدل‌های نااستاندارد}

\begin{explain}
!!a{structure} برای~$\Lang{L_A}$ را استاندارد می‌نامیم اگر با
$\Struct{N}$ همریخت باشد. اگر !!a{structure} با~$\Struct{N}$ همریخت
نباشد، نااستاندارد نامیده می‌شود.
\end{explain}

\begin{defn}
!!{structure}~$\Struct{M}$ برای~$\Lang{L_A}$
\emph{نااستاندارد} است اگر با~$\Struct{N}$ همریخت نباشد.
!!{element}sی چون $x \in \Domain{M}$ که برای عددی $n \in \Nat$ با
$\Value{\num{n}}{M}$ برابرند،
\emph{اعداد استاندارد} (در~$\Struct{M}$) و آن‌هایی که چنین نیستند
\emph{اعداد نااستاندارد} نامیده می‌شوند.
\end{defn}

\begin{explain}
بنا بر
\olref[stm]{prop:standard-domain}، هر !!{structure} استاندارد
برای~$\Lang{L_A}$ تنها !!{element}s استاندارد دارد. در نتیجه، وجود یک
!!{structure} نااستاندارد تضمین می‌کند که !!{element} نااستاندارد
است. عکس این گزاره برای !!{structure} دلخواه برقرار نیست: ممکن است
دامنه تنها از مقادیر عددنماهای استاندارد تشکیل شده باشد ولی تعبیر
نمادها ساختار را نااستاندارد کند. بااین‌حال، اگر $\Struct{M}$ مدلِ
$\Th{Q}$ نیز باشد،
\olref[stm]{prop:thq-standard} نشان می‌دهد که هر $\Struct{M}$
نااستاندارد دست‌کم یک عدد نااستاندارد دارد.
\end{explain}

\begin{prop}
اگر !!a{structure}~$\Struct{M}$ برای~$\Lang{L_A}$ عددی نااستاندارد
داشته باشد، $\Struct{M}$ نااستاندارد است.
\end{prop}

\begin{proof}
خلاف آن را فرض کنید؛ یعنی فرض کنید $\Struct{M}$ استاندارد است اما عدد
نااستاندارد~$x$ را دربر دارد. بگذارید
$g\colon \Nat \to \Domain{M}$ یک همریختی باشد. به‌آسانی (با استقرا
روی~$n$) دیده می‌شود که
$g(\Value{\num{n}}{N}) = \Value{\num{n}}{M}$. به بیان دیگر، $g$
اعداد استاندارد~$\Struct{N}$ را به اعداد استاندارد~$\Struct{M}$
می‌نگارد. اگر $\Struct{M}$ عددی نااستاندارد داشته باشد، $g$ نمی‌تواند
!!{surjective} باشد، برخلاف فرض.
\end{proof}

\begin{prob}
به یاد آورید که $\Th{Q}$ شامل اصل‌های زیر است:
\begin{align*}
& \lforall[x][\lforall[y][(\eq[x'][y'] \lif \eq[x][y])]] \tag{$!Q_1$}\\
& \lforall[x][\eq/[\Obj 0][x']] \tag{$!Q_2$}\\
& \lforall[x][(\eq[x][\Obj 0] \lor \lexists[y][\eq[x][y']])] \tag{$!Q_3$}.
\end{align*}
!!{structure}sی چون $\Struct{M_1}$، $\Struct{M_2}$ و
$\Struct{M_3}$ بیابید چنان‌که
\begin{enumerate}
\item $\Sat{M_1}{!Q_1}$، $\Sat{M_1}{!Q_2}$ و $\Sat/{M_1}{!Q_3}$؛
\item $\Sat{M_2}{!Q_1}$، $\Sat/{M_2}{!Q_2}$ و $\Sat{M_2}{!Q_3}$؛ و
\item $\Sat/{M_3}{!Q_1}$، $\Sat{M_3}{!Q_2}$ و $\Sat{M_3}{!Q_3}$.
\end{enumerate}
آشکار است که کافی است برای هر~$i$، تنها
$\Assign{\Obj{0}}{M_i}$ و $\Assign{\prime}{M_i}$ را مشخص کنید.
\end{prob}

\begin{explain}
مشخص‌کردن !!{structure}s نااستاندارد برای~$\Lang{L_A}$ به‌اندازهٔ کافی
آسان است. برای نمونه، ساختاری با !!{domain}~$\Int$ بگیرید و
همهٔ نمادهای غیرمنطقی را به شیوهٔ معمول تعبیر کنید. چون هیچ عدد منفی
مقدار~$\num{n}$ برای عددی~$n$ نیست، این ساختار نااستاندارد
است. البته این ساختار به معنای صادق‌کردن همان جمله‌هایی که
$\Struct{N}$ صادق می‌کند،
\emph{مدلِ} حساب نخواهد بود. برای نمونه،
$\lforall[x][\eq/[x'][\Obj{0}]]$ کاذب است. بااین‌حال می‌توانیم وجود
مدل‌های نااستاندارد حساب را با استفاده از قضیهٔ فشردگی به‌آسانی ثابت
کنیم.
\end{explain}

\begin{prop}
بگذارید $\Th{TA} = \Setabs{!A}{\Sat{N}{!A}}$ نظریهٔ~$\Struct{N}$
باشد. $\Th{TA}$ !!a{enumerable} مدل نااستاندارد دارد.
\end{prop}

\begin{proof}
زبان~$\Lang{L_A}$ را با !!{constant} تازه~$c$ بسط دهید و مجموعهٔ
!!{sentence}s زیر را در نظر بگیرید:
\[
\Gamma = \Th{TA} \cup \{\eq/[c][\num{0}], \eq/[c][\num{1}],
\eq/[c][\num{2}], \dots\}.
\]
هر مدل~$\Struct{M^c}$ از~$\Gamma$ دارای !!a{element} چون
$x = \Assign{c}{M}$ خواهد بود که نااستاندارد است، زیرا برای همهٔ
$n \in \Nat$، $x \neq \Value{\num{n}}{M}$. همچنین آشکارا
$\Sat{M^c}{\Th{TA}}$، زیرا $\Th{TA} \subseteq \Gamma$. اگر با
فراموش‌کردن~$c$، $\Struct{M^c}$ را به !!a{structure}~$\Struct{M}$
برای~$\Lang{L_A}$ بدل کنیم، دامنهٔ آن هنوز شامل~$x$ نااستاندارد است
و نیز $\Sat{M}{\Th{TA}}$. گزارهٔ اخیر از اینجا تضمین می‌شود که~$c$
در~$\Th{TA}$ رخ نمی‌دهد. پس کافی است نشان دهیم $\Gamma$ مدل دارد.

برای این کار از قضیهٔ فشردگی استفاده می‌کنیم. اگر هر زیرمجموعهٔ متناهی
از~$\Gamma$ ارضاپذیر باشد، خود~$\Gamma$ نیز ارضاپذیر است. هر
$\Gamma_0 \subseteq \Gamma$ متناهی را در نظر بگیرید.
$\Gamma_0$ شامل تعدادی !!{sentence}s از~$\Th{TA}$ و تعدادی از صورت
$\eq/[c][\num{n}]$ است، اما تنها به‌شمار متناهی. اگر هیچ جمله‌ای از
این صورت در~$\Gamma_0$ نیست، $k=0$ را بگیرید؛ وگرنه فرض کنید $k$
بزرگ‌ترین عددی باشد که $\eq/[c][\num{k}] \in \Gamma_0$. با بسط‌دادن
$\Struct{N}$ و قراردادن $\Assign{c}{N_k} = k+1$،
$\Struct{N_k}$ را تعریف کنید. آنگاه $\Sat{N_k}{\Gamma_0}$: اگر
$!A \in \Th{TA}$، چون $\Struct{N_k}$ جز در تعبیر~$c$ از هر جهت
همانند~$\Struct{N}$ است و~$c$ در~$!A$ رخ نمی‌دهد،
$\Sat{N_k}{!A}$. همچنین چون $n \le k$ و
$\Value{c}{N_k} = k+1$،
$\Sat{N_k}{\eq/[c][\num{n}]}$. پس هر زیرمجموعهٔ متناهی~$\Gamma$
ارضاپذیر است و قضیهٔ فشردگی مدلی برای~$\Gamma$ به دست می‌دهد.

از آنجا که زبان و~$\Gamma$ شمارا هستند، قضیهٔ لوونهایم--اسکولم تضمین
می‌کند که~$\Gamma$ مدلی شمارا~$\Struct{M^c}$ دارد. با گرفتن تحویل آن
به~$\Lang{L_A}$، چنان‌که در بالا توضیح داده شد، مدل شمارا و
نااستانداردی از~$\Th{TA}$ به دست می‌آوریم.
\end{proof}

\end{document}
